\section{Evaluation}
\label{sec:eval}

\begin{table}[tb]
\footnotesize
\centering

\caption{Number of injection points, monitor points, and integration tests in each system.}
% \vspace{-3mm}

\begin{tabular}{l|ccccc}
\toprule
\textbf{System} & \textbf{Loop} & \textbf{Exception} & \textbf{Negation} & \textbf{Branch} & \textbf{Test} \\
\midrule
HDFS 2 & 2067 & 2316 & \phantom{0}770 & 26941 & 2674 \\
HDFS 3 & 2736 & 2707 & \phantom{0}974 & 36745 & 4426 \\
HBase & 3227 & 2369 &            1002 & 38192 & 4436 \\
Flink & 2860 & 2619 & \phantom{0}447 & 27947 & 2036 \\
OZone & 1361 & 1395 & \phantom{0}395 & 18212 & 1219 \\
\bottomrule
\end{tabular}
% \vspace{-2mm}
\label{tab:static-result}
\end{table}

We evaluate \tool on the latest versions of \NumOfSysTested popular distributed systems, namely, HDFS 3.4.1~\cite{apache-hdfs3}, HDFS 2.10.2~\cite{apache-hdfs}, HBase 2.6.0~\cite{apache-hbase}, Flink 1.20.0~\cite{apache-flink}, OZone 1.4.0~\cite{apache-ozone}.
The fault injection tests and beam search are performed on two Ubuntu 22.04 servers: one with two Intel Xeon Gold 5220R CPUs and 512 GB of memory, and the other with two AMD EPYC 7313 CPUs and 256 GB of memory. Each profile and injection run is executed in a Docker container with Java 1.8. Each JVM instance is limited to 6 CPU cores and a heap size of 32GB. 



\Cref{tab:static-result} summarizes the number of injection points, monitor points, and test workloads identified in each system. 
In all the evaluations, we use a beam size of 5 million chains. 


\subsection{New \VCs}
\label{sec:eval:new-bug}

\begin{table*}[tb]
\footnotesize
% \setlength{\tabcolsep}{3pt}
\centering

\caption{All \NumOfNewVC new \vcs detected by \tool. In general, ``DN'' stands for DataNode, ``NN'' stands for NameNode, and ``IBR'' stands for incremental block report. ``Cycle'' column shows the types of faults in each cycle. ``Alloc.'' column shows the phase number in 3PA protocol that each bug is detected in. A ``\Yes'' in column ``Rnd.?'' means that the bug is detectable under random allocation of test budgets. ``Alt.?'' column shows whether the bug can be triggered by the strategy in \Cref{sec:eval:alt}. ``JIRA\#'' column lists the bug ID in Apache issue tracker.}
% \vspace{-3mm}
\begin{tabular}{llllccccl}
\toprule

\textbf{System} &
\textbf{\#} &
\textbf{Delayed Task} & 
\textbf{Other Faults} &
\textbf{Cycle*} & 
\textbf{Alloc.} & 
\textbf{Rnd.?} & 
\textbf{Alt.?} &
\textbf{JIRA\#}\\
\midrule 

HDFS 2 & 
1) & 
Lease recovery &
IOE in IBR; IOE in write pipeline & 
1D | 2E | 0N & 
1 &
 &
\Yes &
17661\\

 & 
2) & 
Edit log flushing &
IOE in IBR after NN failover & 
1D | 1E | 0N & 
1 &
\Yes &
 &
17836 \\

 & 
3) & 
Block recovery &
IOE in block recovery & 
1D | 1E | 0N & 
1 &
\Yes &
\Yes &
17662\\

 & 
4) & 
Write pipeline &
IOE in block recovery, IBR, and write pipeline & 
1D | 3E | 0N & 
2 &
 &
 &
17837 \\

 & 
5) & 
Block cache &
IOE in write pipeline; DN timeout & 
1D | 1E | 1N & 
2 &
 &
 &
17660 \\ 

 & 
6) & 
IBR &
IOE in IBR & 
1D | 1E | 0N & 
3 &
 &
 &
17780 \\
\midrule 

HDFS 3 & 
1) & 
Block deletion &
IOE in write pipeline; DN timeout & 
1D | 1E | 1N & 
2 &
\Yes &
 &
17838 \\ 
 
 & 
2) & 
Block reconstruction; IBR &
DN timeout; IOE in replication & 
2D | 1E | 1N & 
3 &
 &
 &
17782 \\ 
\midrule 

HBase & 
1) & 
Write ahead log (WAL) &
Premature \texttt{EndOfFile} in WAL & 
1D | 0E | 1N & 
1 &
\Yes &
\Yes &
29600\\

 & 
2) & 
Region assignment &
IOE in assignment RPC; Node exclusion & 
1D | 1E | 1N & 
3 &
 &
 &
29006 \\ 
\midrule 

Flink & 
1) & 
Task worker &
Head task failure; Sink task cancellation & 
1D | 2E | 0N & 
1 &
\Yes &
 &
38367 \\

 & 
2) & 
Aggregation task &
Task state transition failure; Barrier task failure & 
1D | 2E | 0N & 
2 &
\Yes &
 &
38368 \\
\midrule

OZone & 
1) & 
Container report queue &
Event queue dispatch failure & 
1D | 0E | 1N & 
1 &
 &
 &
13020\\

 & 
2) & 
Heartbeat handling &
IOE in pipeline construction; Pipeline unhealthy & 
1D | 1E | 1N & 
2 &
\Yes &
\Yes &
$11856_1$\\

 & 
3) & 
Replication command handling &
IOE in replication and pipeline construction & 
1D | 2E | 0N & 
3 &
 &
 &
$11856_2$ \\
 
\bottomrule

\end{tabular}

\hspace{48em}\scriptsize{*D: Delay; E: Exception; N: Negation}


\label{tab:bug-detected}
\end{table*}



\Cref{tab:bug-detected} shows the \NumOfNewVC new \vcs detected by \tool in all \NumOfSysTested distributed systems, many of which involve retries and the failure recovery logic. This aligns with the findings of \citet{qian2023vicious} on \vcs.  
\NUMOfConfirmedVC of them have been confirmed, with \NumOfFixedVC fixed. 
Two bugs detected in HDFS 2 are also detected in HDFS 3, we do not list them in \Cref{tab:bug-detected} due to duplication. 

\Cref{tab:bug-detected} also details the characteristics of the fault injection runs for each bug. Column ``Cycle'' shows the number of faults injected of each type.
Column ``Alloc.'' shows the test budget allocation phase after which all the causal relationships of each bug are discovered. All three phases of the 3PA protocol contributes to detecting new bugs. 

To further demonstrate 3PA protocol's effectiveness, we compare it with a random allocation protocol, running the same number of fault injection runs as the 3PA protocol with randomly selected (fault, test) combinations. A ``\Yes'' in column ``Rnd.?'' of \Cref{tab:bug-detected}  means that the bug can be detected using random allocation. The random selection allocation protocol, albeit effective in some systems (i.e., Flink), generally under-performs the 3PA protocol in \tool. 

Most new failures detected by \tool require only one delay injection, with varying numbers of exception and negation injections. 
This is in line with the bug dataset provided by \citet{qian2023vicious}, where majority of the cases only requires a single contention to be triggered. 

\Vcs reported by \tool that involve multiple delay injections generally contribute to the false positive cases. 
However, most of the false positive cases 
are valid fault propagations, 
but categorized as false positive 
because they are known (and accepted by the developer) 
contentions between operations in the system 
(e.g., contention between the HDFS clients performing read and write). 
Details will be discussed in the \Cref{sec:eval:fp}.


\subsection{Comparison with Alternative Strategy}
\label{sec:eval:alt}

We compare \tool with a naive strategy that 
injects a single fault into a system and monitors whether it causes itself 
(e.g., delays a single loop in the system and monitors if its iterations increase). 
This is to evaluate whether the triggering conditions of each detected bug span multiple tests.

Column ``Alt.?'' in \Cref{tab:bug-detected} shows that about \NoNaivePctg (\NumOfNewVCNaive out of \NumOfNewVC) of the \vcs detected by \tool cannot be triggered by the naive strategy. 
This underscores the effectiveness of \tool's causal stitching, the complexity needed to trigger these failures, and the limitations of single-fault injection tools. Even when the naive strategy works, \tool offers detailed insights into the cycle's behavior.

\subsubsection{Comparison with Existing Fuzzing Techniques}

In addition to the naive strategy above, we compare \tool with blackbox fuzzing techniques. Specifically, we compare with Jepsen~\cite{jepsenio} and its Python reimplementation, Blockade~\cite{blockadetest}. Jepsen is a Clojure-based blackbox distributed system fuzzer, widely being used in the security research community~\cite{meng2023greybox,zou2025blackbox} for distributed system testing.

We apply Jepsen on Flink~\cite{jepsenflink} and Blockade on OZone~\cite{blockaeozone} with their existing test cases. Due to the lack of Jepsen or Blockade integration and test workload, HBase and both versions of HDFS are not included in this comparison.

The result shows that none of the 
\vcs detected by \tool can be detected by fuzzing-based techniques. This highlights the necessity of causal stitching and fine-grained fault injection in \tool for detecting \vcs. 

\subsection{Case Study}
\label{sec:eval:case-study}

In this section, we present case studies on two of the confirmed \vcs.

\subsubsection{Region Deployment Retry in HBase}

\begin{figure}
    \centering
    \includegraphics[width=0.8\columnwidth]{figure/motivating_example-illus-new.pdf}
    % \vspace{-3mm}
    \caption{The case study. ``RS'' stands for RegionServer, and ``IOE'' stands for IOException. The yellow box indicates a delay.}
    \label{fig:motivating-exp}
    % \vspace{-3mm}
\end{figure}


The second HBase bug in \Cref{tab:bug-detected} is a \vc 
similar to the AWS motivating example in \Cref{sec:intro}. 
HBase is a distributed database system 
that divides tables into regions managed by multiple RegionServers (RS) under master nodes.

At high level, RSes in a heavily loaded cluster may throw IOExceptions (IOEs) when handling RPC requests. The \vc happens in a cluster with a lot of write and table creation requests, which place significant load on the RSes managing these tables (\circledtext{1} in \Cref{fig:motivating-exp}). Some region assignment operations on these RSes may time out and throw IOEs (\circledtext{2}). 
When an RS throws an IOE, it is excluded from the master node's \texttt{Favored\-Stochastic\-Ba\-lan\-cer}, which requires at least three live RSes to function properly. A reduced number of live RSes causes the load balancer to fail (\circledtext{3}). 
An improper handling logic blindly retries the assignment indefinitely (\circledtext{4}), further increasing the load on the RSes (\circledtext{5}). This creates a self-sustaining cycle, leaving the cluster unable to process additional region assignment requests.


To trigger this bug in a testing environment with a single-fault (i.e., delay) injection, we need a workload with the following conditions:
\begin{enumerate}[leftmargin=5ex, noitemsep, topsep=0pt]
\item It contains many region assignment requests. 

\item A cluster configuration prone to be overloaded (3 nodes). 

\item Load balancer configured to \texttt{F.S.Balancer}. 

\item A workload long enough to observe the cycle. 
\end{enumerate}

However, no single test in HDFS satisfies 
all these conditions (as shown in \Cref{sec:eval:alt}), 
and 
\tool detects this cycle by injecting one delay, one exception, and one return-value negation into 
\textbf{three separate tests} $t_1$, $t_2$, and $t_3$. 

First, the delay is injected into the region deployment loop in 
$t_1$
with multiple table creation and clones (cf. \circledtext{1} in \Cref{fig:motivating-exp}). This overloads the cluster, causing IOEs in the region assignment request handler (cf. \circledtext{2}). 
This causal relationship is only detected in 
$t_1$, 
because other tests only exercise limited number of region assignments. 

In the second experiment, 
the same IOE is injected into the region assignment request handler 
in a test ($t_2$) for RS fault tolerance, 
triggering a negated return value in the load balancer's status checker (cf. \circledtext{3}). 
This causal relationship is only detected 
in $t_2$, 
because $t_2$ is the only test that 
uses 
\texttt{Favored\-Stochastic\-Balancer} 
with a cluster of 3 nodes, 
and 
the error detector \texttt{canPlaceFavoredNodes()} 
in \texttt{Favored\-Stochastic\-Balancer} 
only runs into an error (negation) 
if the cluster has fewer than 3 nodes 
(e.g., $t_3$ uses \texttt{Favored\-Stochastic\-Balancer} but has 5 nodes). 


Finally, in the third experiment ($t_3$), the return-value negation is injected into the load balancer during a test for the balancer itself, increasing the iterations of the region deployment loop (cf. \circledtext{4} and \circledtext{5}), indicating an increase in the workload being handled that can potentially cause delays. 
This causal relationship is only observed in 
$t_3$, 
because 
it is the only test that uses 
\texttt{Favored\-Stochastic\-Balancer} configuration 
and has a long enough workload to observe the increased iterations 
(not in $t_2$ because it exits prematurely after the IOE). 

\subsubsection{Bypassed IBR Throttling in HDFS}

The sixth bug detected on HDFS 2 is a logic error that a failed 
incremental block report (IBR) is retried immediately at the next heartbeat, 
incorrectly ignoring the configured IBR interval. 
The unconstrained, immediate retry of IBR 
can cause a \vc if the original IBR failure is 
due to a server overload. 
This happens when the cluster 
is under high user traffic with many modifications, 
which triggers many large-sized IBRs.

\tool detects this bug by linking two injections in two different tests $t_1$ and $t_2$: 1) injecting an IBR processing delay into a workload with over 5,000 blocks to test the load balancer, and 2) injecting an RPC exception for IBR into a workload testing the IBR report interval configuration. 

In fault injection using $t_1$, the RPC exception is observed, but not an increase in IBR count. 
Because $t_2$ is the only one with IBR throttling configured. 
Without it, IBRs are sent with every heartbeat. 
As a result, 
although IBR failures occur in the first workload, they are still sent at the original frequency and cannot be detected by execution trace comparison. 

In fault injection using $t_2$, 
the IBR increment is observed 
after the RPC exception, indicating a potential delay. 
However, $t_2$ involves only 8 file blocks 
compared to 5,000 in $t_1$, 
making IBR processing delays less likely to cause timeouts.


\subsection{Cycle Clustering and False Positive Rate}
\label{sec:eval:fp}

\begin{table}[tbp]
\centering
\footnotesize
\setlength{\tabcolsep}{3pt}

\caption{Number of \vcs reported by \tool and the clustering results. Column ``System'' is the targeted distributed system. Columns ``Cycle'', ``Cluster'', and ``TP'' shows the number of cycles reported, distinct cycle clusters, and true positive clusters respectively. The numbers outside the parentheses are from a beam search with unlimited number of injected faults, while the numbers inside are from a beam search limiting to one delay injection.}
% \vspace{-3mm}
\begin{tabular}{lccc||lccc}
\toprule
\textbf{System} & \textbf{Cycle} & \textbf{Cluster} & \textbf{TP} &
\textbf{System} & \textbf{Cycle} & \textbf{Cluster} & \textbf{TP}\\
\midrule 

HDFS 2 & \phantom{0}38 (23) & 15 \phantom{0}(9) & 6 (6) &

Flink & 48 (27) & 35 (17) & \phantom{0}2 (2)\phantom{x} \\

HDFS 3* & 149 (59) & 36 (14) & 4 (3) &

OZone & 29 (17) & 11 \phantom{0}(7) & \phantom{0}3 (3)\phantom{x} \\

HBase & \phantom{0}72 (26) & 24 \phantom{0}(7) & 2 (2) &

\textbf{Total}* & & \multicolumn{2}{r}{\textbf{\NumOfNewVCAllowDupe (\NumOfNewVCOneDelayAllowDupe)}} \\

\bottomrule
\end{tabular}

\hspace{9em}\scriptsize{*Including two duplicated clusters also detected in HDFS 2}



\label{tab:cluster-result}
\end{table}

\Cref{tab:cluster-result} shows the number of \vcs, distinct failure clusters, and the true positive reported by \tool on each system. 
The numbers outside the parentheses are from a beam search without an upper bound on the number of injected faults, while the numbers inside the parentheses are from a beam search for cycles with at most one delay injection but unlimited exceptions or negations.

Limiting the beam search to one delay injection generally reduces false positives while still identifying most failures. Users can adjust the number of injected faults to balance between accuracy and completeness.



\subsubsection{Effectiveness of Cycle Clustering}

We manually inspect all the \vcs and the \textit{causally equivalent fault} clusters on HDFS 2 and confirm that the clustering algorithm effectively groups similar cycles together. Examples of the clustered cycles involve faults from 1) the write pipeline, 2) DataNode command generation and processing. 

However, we do observe false positives in the clustering results due to 1) non-determinism and 2) insufficient test execution. For example, a fault cluster in a HDFS run mixed four correctly-clustered data-race faults with two false positives: a failed security check whose unique causal consequence would have been revealed with more test cases, and a logging error clustered due to non-determinism. This could cause missed \vcs in detection if the false positive fault is selected during the 3PA allocation, though their impact is limited as long as the majority are clustered correctly.

The higher numbers of cycles and clusters reported in HDFS 3 are due to the extensive usage of asynchronous tasks and event queues, 
where 
errors in the issuer 
and the executor of asynchronous tasks 
are processed by different error handlers. 
Compared to synchronous execution, 
this 
increases the number of error handlers, 
resulting in 
a larger number of 
cycles and fault clusters.

\subsubsection{False Positive Analysis}

The false positive clusters reported by \tool are mainly due to the following reasons:
\begin{enumerate}[wide, labelwidth=!, labelindent=\parindent, noitemsep, topsep=0pt]
\item 
In cycle searches with unlimited fault injections, many false positives (about 70\% on HDFS 2) are due to ``expected'' contentions, such as between HDFS clients with heavy read and write operations. These cases require solutions such as throttling or capacity increases. 
We do not consider them as true positives. 

\item 
Increased loop iterations don't always indicate delays, especially if throttling mechanisms are already in place. These mechanisms are often implemented ad-hoc, making them difficult for static analyzers to identify.

\end{enumerate}

From our experience, 
the false positives 
due to the above reasons can be 
filtered out manually 
relatively easily, 
given adequate understanding of the system under test. 

Theoretically, \tool could still report a \vc involving fault propagation chains from tests with contradictory conditions after our local compatibility check. However, we do not observe such cases in the clusters reported by \tool.

\subsection{Performance Overhead}
\label{sec:eval:perf}

\tool's instrumentation and monitoring introduces an average of 185\% runtime overhead in the profile run, ranging from 63\% to 376\%. This overhead is primarily caused by branch tracing and call stack recording, 
and can be further reduced using hardware-based tracing such as IntelPT~\cite{henschelintelpt}. 
While sampling can be used to reduce the runtime overhead, we opt not to implement sampling in branch tracing to maximize the detection accuracy as \tool is intended to be deployed in the testing environment instead of the production environment.

\subsection{Effort of Applying \tool to a New System}

\tool is designed to be system-agonistic and requires as little input from users as possible. Applying \tool to a new system only requires two steps: 1) compiling the target system, and 2) setting up the test environment with system dependencies. Two graduate students who have no knowledge of the tool can independently apply \tool to a new system within half a day without consulting the authors. 